\chapter{Núcleo del trabajo}

\section{Descripción general de la solución}

La solución actual se estructura como un sistema de agentes autónomos entrenados dentro del simulador. Cada vehículo dispone de un controlador de IA que percibe el entorno, ejecuta una red neuronal, aplica acciones al sistema físico y registra métricas de comportamiento. Por encima de los vehículos existe un gestor evolutivo que crea poblaciones, asigna pesos, evalúa resultados, selecciona los mejores individuos y guarda modelos.

El sistema se completa con elementos de tráfico que comunican contexto a los controladores: semáforos, señales de stop, señales de ceda y zonas de cruce. Estos elementos no son meros objetos visuales, sino actores con lógica propia que activan zonas de detección, informan al vehículo de líneas de parada, controlan estados y aplican penalizaciones o bonus según el comportamiento observado.

\section{Subsistemas}

\subsection{Gestor evolutivo}

\texttt{BP\_EvolutionManager} coordina el ciclo de entrenamiento y las ejecuciones de test. Al comenzar comprueba si existe un modelo histórico en \texttt{SuperCarModel}. Si lo encuentra, carga sus pesos; en caso contrario, inicializa una red aleatoria que servirá como punto de partida. Después actualiza la lista de puntos de aparición permitidos para el modo activo y crea la población de vehículos, asignando a cada coche su controlador y su red.

Durante el entrenamiento conviven individuos procedentes del mejor modelo histórico, del mejor modelo obtenido en la sesión y de inicializaciones aleatorias. Los primeros se conservan o se mutan para explotar soluciones que ya han funcionado, mientras que los aleatorios mantienen capacidad de exploración. La generación termina cuando todos los vehículos han finalizado o han sido eliminados por alguna condición de muerte. En ese momento el gestor ordena los resultados, actualiza los modelos persistidos cuando corresponde, destruye la población anterior e inicia la siguiente generación.

El modo de test reutiliza la misma infraestructura, pero desactiva la mutación y carga un modelo fijo. También emplea un conjunto de puntos de aparición distinto del utilizado para entrenar. Esta separación evita que una mejora aparente proceda únicamente de evaluar el coche en los recorridos que ya han condicionado su aprendizaje. Al terminar cada bloque, el gestor exporta tanto el resumen de la ejecución como los registros individuales necesarios para el análisis.

Una decisión relevante es la normalización por punto de aparición. La ciudad contiene zonas de dificultad desigual y comparar directamente el fitness de dos spawns puede favorecer al que comienza en una ruta sencilla. Por ello se conserva el mejor valor observado en cada punto y el resultado de un coche se relaciona con esa referencia local antes de seleccionar sus pesos. La normalización no elimina las diferencias del mapa, pero reduce su influencia en la elección del individuo que pasa a la siguiente generación.

\subsection{Controlador de IA y red neuronal}

\texttt{BP\_AiController} actúa como intermediario entre red neuronal y vehículo físico. En cada tick comprueba que el vehículo controlado existe y no está muerto, prepara las entradas de red, ejecuta la inferencia y aplica las salidas al componente de movimiento.

Para evitar saltos bruscos, el controlador dispone de variables de suavizado como \texttt{MaxSteeringRate} y \texttt{MaxThrottleRate}. La versión documentada mantiene un seguimiento estable entre objetivo y entrada aplicada, con gap medio cero en los informes finales, lo que indica que la señal de control se está transmitiendo sin divergencias instrumentales.

La inferencia se calcula manualmente sobre arrays de pesos:

\begin{equation}
\label{eq:capa-oculta}
h_j = \tanh\left(\sum_i w^{IH}_{j,i} x_i\right)
\end{equation}

\begin{equation}
\label{eq:capa-salida}
y_k = \tanh\left(\sum_j w^{HO}_{k,j} h_j\right)
\end{equation}

La ecuación~\ref{eq:capa-oculta} calcula las activaciones ocultas y la ecuación~\ref{eq:capa-salida} obtiene las dos salidas de control, donde \(x_i\) son las entradas normalizadas, \(h_j\) las neuronas ocultas e \(y_k\) las salidas. Los superíndices identifican los dos bloques de pesos de la red. \(IH\) significa \textit{Input--Hidden} y corresponde a las conexiones entre la entrada y la capa oculta. \(HO\), denominado \texttt{WeightsHO} en los Blueprints y anotado en ocasiones como \(OH\), corresponde a las conexiones \textit{Hidden--Output} entre la capa oculta y las dos salidas.

\Needspace{10\baselineskip}
Los pesos se almacenan en arrays unidimensionales. Por esa razón, el número de entradas esperado se obtiene dividiendo la longitud de \texttt{WeightsIH} entre el número de neuronas ocultas. A continuación se comprueba que ese resultado coincide con el tamaño del vector construido por el controlador, que \texttt{WeightsIH} contiene exactamente \(n_{entradas}\times n_{ocultas}\) valores y que \texttt{WeightsHO} contiene \(n_{ocultas}\times n_{salidas}\). Este cálculo sustituye a una comparación rígida con 22 y permite detectar un modelo antiguo aunque cambie la arquitectura. Si faltan pesos o las dimensiones no son coherentes, se activa una guarda; después de cinco errores consecutivos, el coche se elimina con la razón de muerte \texttt{INVALID\_BRAIN} para evitar que continúe circulando con una red inválida.

\subsection{Percepción mediante sensores}

El controlador lanza once trazas en abanico sobre 180 grados. Cada sensor devuelve una distancia normalizada hasta el primer elemento relevante. Los bordes de intersección se tratan con lógica específica: un borde virtual solo interviene en la lectura cuando su intención de dirección y su trigger coinciden con la decisión vigente del vehículo. Los bordes asociados a otras trayectorias se ignoran, de modo que una misma intersección puede contener límites sensoriales distintos para cada maniobra.

\Needspace{6\baselineskip}
La decisión se modela mediante la enumeración \texttt{E\_DirectionIntent}, con hasta cuatro alternativas de salida. Al entrar en un \texttt{BP\_IntersectionTrigger}, el vehículo selecciona aleatoriamente una de las opciones habilitadas para el modo de entrenamiento o de test. El controlador conserva tanto la dirección elegida como el trigger que la originó. Cada \texttt{BP\_IntersectionBorders} declara la pareja dirección--trigger a la que pertenece, lo que permite que los rayos perciban únicamente los límites que delimitan la trayectoria seleccionada.

\Needspace{6\baselineskip}
Las entradas sensoriales se combinan con información de tráfico. Si el vehículo entra en la zona de un semáforo, stop o ceda, el actor correspondiente comunica al controlador el tipo de control, la línea de parada y la distancia de entrada. Esta mezcla de percepción geométrica y contexto normativo permite que la red no dependa únicamente de distancia a obstáculos.

\Needspace{10\baselineskip}
\subsection{Sistema de fitness}

La función de fitness se ha convertido en uno de los elementos centrales del proyecto. Las primeras versiones premiaban principalmente el avance, pero ese criterio por sí solo permitía comportamientos poco útiles: circular demasiado rápido cerca de un obstáculo, corregir hacia el lado equivocado o sobrevivir sin progresar. La función actual combina el avance con la velocidad útil y el centrado estimado a partir de los sensores laterales. También valora que la velocidad sea coherente con la distancia frontal disponible, de modo que aproximarse a un obstáculo sin reducir la marcha deje de ser ventajoso.

\Needspace{8\baselineskip}
Las penalizaciones se aplican cuando el vehículo marcha hacia atrás, permanece atascado, abandona la calzada o colisiona. Hay comprobaciones específicas para las correcciones de dirección y para el incumplimiento de semáforos, stops, cedas y zonas no navegables. Completar correctamente un stop o un ceda produce una bonificación, pero esa recompensa se registra por separado para poder comprobar que no compensa una infracción grave ocurrida en otro momento de la ejecución.

Además del valor final, cada coche acumula métricas de aceleración, frenado, espera, dirección y contexto de parada. Esta instrumentación permite reconstruir qué componentes han dominado el resultado y distinguir, por ejemplo, una finalización por tiempo con recorrido útil de un vehículo inmóvil ante una señal. El desglose fue decisivo después de la demostración fallida del 5 de marzo de 2026, cuando la observación visual no bastó para localizar si el problema procedía de los sensores, de las restricciones o del propio cálculo de fitness.

\subsection{Semáforos}

El sistema de semáforos se compone de un controlador de intersección y de actores de semáforo. El controlador alterna fases norte-sur y este-oeste, incluyendo verde, ámbar y todo rojo. Cada semáforo mantiene materiales dinámicos para representar estado visual y una zona de sensor que detecta vehículos.

Cuando un vehículo entra en la zona, el semáforo comunica estado, línea de parada y distancia. Si el semáforo está en rojo o ámbar restrictivo, el controlador activa contexto de parada. Al salir de la zona, se comprueba si el vehículo atravesó la línea a una velocidad legal. El sistema aprende umbrales de velocidad legal a partir de muestras previas, evitando depender siempre de un valor fijo.

\Needspace{10\baselineskip}
\subsection{Stop y ceda el paso}

Las señales \texttt{BP\_Sign\_Stop} y \texttt{BP\_Sign\_Yield} amplían la interacción normativa. En stop se exige detención suficiente dentro de una ventana dinámica antes de la línea, espera mínima y comprobación de cruce libre. Si se cumple, se aplica bonus y se libera el vehículo; si no, se aplica penalización y muerte por \texttt{STOP\_VIOLATION}. En ceda se monitoriza velocidad, aproximación y disponibilidad de la zona de cruce, con penalización por \texttt{YIELD\_VIOLATION} y bonus al completar correctamente.

La diferencia entre stop y ceda es metodológicamente importante. El stop requiere detención explícita; el ceda permite continuar si la velocidad y la prioridad son compatibles con la seguridad del cruce.

\section{Pipeline experimental}

Cada ejecución genera un registro por vehículo y un resumen por generación. Una vez terminada la sesión, los scripts de Python validan que ambos CSV contienen el número de filas esperado y que sus identificadores pueden alinearse. Solo después de esa comprobación se calculan los indicadores y se generan las gráficas. De esta manera, una fila incompleta o una generación duplicada no queda oculta dentro de una media agregada.

El diagnóstico comienza con \textit{BestFit}, \textit{MeanFit} y tiempo de vida, pero no termina en esas tres cifras. Las razones de muerte se agrupan por familias para separar colisiones, finalizaciones por tiempo y fallos normativos; los bonus de stop y ceda se contrastan con el contexto de frenado; y los resultados se desglosan por punto de aparición. También se revisan las correlaciones entre los componentes del fitness y la estabilidad con la que las consignas de dirección y aceleración llegan al vehículo.

Este procedimiento ha permitido utilizar los datos para decidir el siguiente cambio de implementación. Cuando un indicador empeora, se vuelve al CSV individual para localizar actores, intersecciones o spawns concretos; después se modifica el Blueprint correspondiente y se repite el entrenamiento o el bloque de test. El pipeline cierra así el ciclo entre implementación, simulación y análisis sin depender únicamente de lo observado en pantalla.
